% Options for packages loaded elsewhere
\PassOptionsToPackage{unicode}{hyperref}
\PassOptionsToPackage{hyphens}{url}
\documentclass[
  11pt,
]{article}
\usepackage{xcolor}
\usepackage[margin=1in]{geometry}
\usepackage{amsmath,amssymb}
\setcounter{secnumdepth}{-\maxdimen} % remove section numbering
\usepackage{iftex}
\ifPDFTeX
  \usepackage[T1]{fontenc}
  \usepackage[utf8]{inputenc}
  \usepackage{textcomp} % provide euro and other symbols
\else % if luatex or xetex
  \usepackage{unicode-math} % this also loads fontspec
  \defaultfontfeatures{Scale=MatchLowercase}
  \defaultfontfeatures[\rmfamily]{Ligatures=TeX,Scale=1}
\fi
\usepackage{lmodern}
\ifPDFTeX\else
  % xetex/luatex font selection
  \setmainfont[]{Times New Roman}
\fi
% Use upquote if available, for straight quotes in verbatim environments
\IfFileExists{upquote.sty}{\usepackage{upquote}}{}
\IfFileExists{microtype.sty}{% use microtype if available
  \usepackage[]{microtype}
  \UseMicrotypeSet[protrusion]{basicmath} % disable protrusion for tt fonts
}{}
\makeatletter
\@ifundefined{KOMAClassName}{% if non-KOMA class
  \IfFileExists{parskip.sty}{%
    \usepackage{parskip}
  }{% else
    \setlength{\parindent}{0pt}
    \setlength{\parskip}{6pt plus 2pt minus 1pt}}
}{% if KOMA class
  \KOMAoptions{parskip=half}}
\makeatother
\setlength{\emergencystretch}{3em} % prevent overfull lines
\providecommand{\tightlist}{%
  \setlength{\itemsep}{0pt}\setlength{\parskip}{0pt}}
\usepackage{bookmark}
\IfFileExists{xurl.sty}{\usepackage{xurl}}{} % add URL line breaks if available
\urlstyle{same}
\hypersetup{
  hidelinks,
  pdfcreator={LaTeX via pandoc}}

\author{}
\date{}

\begin{document}

\section{Response to the Academic
Editor}\label{response-to-the-academic-editor}

\textbf{Manuscript:} \emph{Proxy-Based Diagnostics of Quantum Oracle
Sketching Robustness for Non-IID Sensor and Telemetry Streams}\\
\textbf{Journal:} \emph{Sensors} (MDPI) · \textbf{Manuscript ID:}
sensors-4470240

\textbf{Editor's comment.} \emph{``I strongly invite the Authors to not
include preprints among references. These works are still undergoing the
peer-review process, thus meaning that their content may change prior to
publication (if they will be finally published). The Authors can
substitute them with related works already published.''}

\textbf{Response.} We thank the Academic Editor and agree with the
principle. We audited all 54 references: every entry was checked against
the publisher's own metadata, and all 51 DOIs resolved and matched the
cited title, first author, venue, and year. The audit found two
preprints. One had in fact been published and has been replaced: the
Kallaugher, Parekh and Voronova paper now carries its published version
(SOSA 2025, pp.~9--45, DOI 10.1137/1.9781611978315.2). The same sweep
also fixed a transposed article number in the Su, Guo and Wang (2024)
entry, which now reads \emph{Information Sciences} \textbf{676}, article
120799, DOI 10.1016/j.ins.2024.120799.

One preprint remains: Zhao et al., ``Exponential quantum advantage in
processing massive classical data'' (arXiv:2604.07639). On 16 August
2026 we re-checked whether it has been published or accepted anywhere.
It has not: the arXiv record is unchanged, no bibliographic record we
identified indicates a published or in-press version, and the authors'
own publication pages still list only the arXiv posting. We retain it
because it is not cited as evidentiary support for a claim of ours; it
is the specific framework under study. The paper examines how that
framework's correlation assumptions behave on non-IID streams; citing a
related published work in its place would attribute its theorems to
authors who did not prove them, and removing it would leave the
framework the paper analyses unattributed.

We have, however, reduced the manuscript's reliance on it to the minimum
necessary for accurate attribution. Its in-text citations were cut from
eight to six by consolidating two repeats. The entry is labelled
``Preprint'' and is formatted in accordance with the MDPI Chicago
reference style for unpublished manuscripts and working papers, and a
footnote at its first citation states this status and explains why the
preprint is not used as evidentiary support for the paper's empirical
conclusions: no quantum algorithm is implemented or simulated, and every
inherited quantity is listed with its provenance and epistemic status in
Table 1. We also added a strong peer-reviewed published anchor of the
same general kind --- Kallaugher, Parekh and Voronova (STOC 2024,
pp.~1805--1815, DOI 10.1145/3618260.3649709), which establishes an
exponential quantum--classical space separation for a natural streaming
problem --- co-cited with it in the Related Work and in the footnote. In
doing so, we corrected a Related Work sentence that had credited the
2021 FOCS paper with an exponential separation; the FOCS 2021 result
provides a polynomial quantum advantage, whereas the STOC 2024 work
establishes the exponential separation. The revised bibliography now
contains 55 entries, including this newly added published anchor, and we
re-checked that every in-text citation matches exactly one bibliography
entry.

Should the work appear in a peer-reviewed venue before this manuscript
reaches proofs, we will substitute the published version at that stage.
If the Editor considers the presence of even this single, explicitly
labelled preprint unacceptable, we would be grateful for the Editor's
preferred remedy and will follow it; we have not removed the attribution
unilaterally because doing so would leave the specific framework under
study without an appropriate bibliographic source.

On behalf of all authors,

\textbf{AbdelMoniem Helmy} (corresponding author)
\href{mailto:abdelmoniem.hafez@cu.edu.eg}{\nolinkurl{abdelmoniem.hafez@cu.edu.eg}}
· ORCID 0000-0001-5996-6019

\end{document}
